\exnumber{04}
\extitle{UX V - Ex 5 - Wireframes}
\exfiles{ex04/fichier Figma, nomme Name\_FirstName\_DeliverableName\_Date\_VersionNumber}
\exauthorize{Those named in the exercise}
\makeheaderfiles

\textbf{Instructions :}

Pick a product (a record, a song or an album) from your music label website.

Based on the new taxonomy (content categorization), show the different paths a user has to take to reach the product with a series of wireframes.

\textbf{Pistes :}

\begin{itemize}
\item Drawing a user journey first might be helpful before rushing into wireframes
\item Drawing paper prototypes before going digital might save you some time ;)
\item Several paths are encouraged
\end{itemize}

Vocabulary

\begin{itemize}
\item Basse fidélité : prototypes papier - exemple \url{https://miro.medium.com/max/1400/1*5nUfqqA2gjdAYHagjbHA5w.jpeg}
\item Moyenne fidélité : wireframes, noir et blanc, seulement les éléments principaux, sur un outil de prototype - exemple \url{https://miro.medium.com/max/1400/1*eGSfaxwYbxUFomYbyO6GWw.png}
\item Haute fidélité : wireframes avec couleur, styles, détails graphiques et micro-interactions - exemple \url{https://miro.medium.com/max/1400/1*Xn0HSKAvhr4TZzC9lN5udw.gif}
\end{itemize}

\textbf{Liste de sites :}

\begin{itemize}
\item Universal Music Latin America \url{https://www.universalmusica.com/}
\item Delta Records \url{https://deltarecords.net/}
\item Zoho Music \url{https://www.zohomusic.com/}
\item Luaka Bop \url{https://www.luakabop.com/}
\item Rimas Music \url{https://rimasmusic.com/}
\item D'autres... tant que c'est un label de musique sud-américain ;-)
\end{itemize}

\textbf{Livrables :}

\begin{itemize}
\item Wireframes (6 to 12 screens) from the home page to the product page, on Figma
\item Déposez le fichier Figma sur Github, nommé ainsi : « Name\_FirstName\_DeliverableName\_Date\_VersionNumber ». Exemple : « Doe\_John\_MidFiPrototype\_05242024\_V1 ».
\item Ajoutez un titre dans le fichier.
\item Soignez l'aspect global du fichier. Il doit être clair, simple et lisible. Vous pouvez regarder canva.com, sans surcharger le dessin de détails.
\end{itemize}

\textbf{Ressources :}

\begin{itemize}
\item How to start sketching and wireframing \url{https://uxdesign.cc/how-to-start-sketching-and-wireframing-84a821f092e2}
\item Why should product teams use wireframes more often? \url{https://uxdesign.cc/why-should-product-teams-use-wireframes-more-often-60e34a2bc55}
\item The 5-pass reduction wireframing: A minimalist UX technique \url{https://uxdesign.cc/the-5-pass-reduction-wireframing-a-minimalist-ux-technique-aa415aad9ce2?sk=a218ac4db9d1049249bd381fbc5edeee}
\item The art of designing good wireframes \url{https://uxdesign.cc/the-art-of-designing-good-wireframes-8a2a9c10a95b}
\end{itemize}

\textbf{Outils recommandés :}

\begin{itemize}
\item Figma
\end{itemize}

